% Root and D completed independent literal audits before input.
\subsection{Projective opposite corners of a strong pair}

\begin{proposition}\label{C:finite-target-strong-pair-twenty}
Let $S$ be a finite strongly connected two-object category and
$F:S^{\mathrm{op}}\to\mathrm{Set}$ a finite right act with both
stalks nonempty. Suppose its act category has at most twenty
arrows. Its classifying space cannot satisfy all four conditions
of Theorem~\ref{C:finite-target-small-act-twenty}: simple
connectivity and rational projective-plane homology, an integral
degree-two generator with nonzero primitive square modulo
torsion, nullity of that generator on every finite group, and
Betti numbers $(1,0,1)$ in degrees two through four over
$\mathbb F_{65521}$.
\end{proposition}
\begin{proof}
Suppose all four conditions hold. Write $s=|S|$ and
$t=|F(a)|+|F(b)|$, so $s+t+1\leq20$ and $t\geq2$.
If one object of $S$ is a retract of the other, delete it
from the entire category. This is an actual homotopy equivalence
to an ordinary act cone with fewer arrows. All four conditions
survive that equivalence, contradicting
Theorem~\ref{C:finite-target-small-act-twenty}.
Assume no such retract exists.

If $s\geq16$, the two stalks have total size at most three.
Lemma~\ref{C:two-states-objectwise-act-cups} gives zero rational
products, a contradiction. We may assume $s\leq15$.
If $H^2(BS;\mathbb Q)=0$, all degree-two classes of the act
cone are relative to $BS$, and the suspension quotient gives
the same contradiction. Otherwise total minimal row and column
counts are both at least three. The two corner identities are
outside the minimal ideal. A nontrivial minimal group would
therefore require at least twenty arrows in $S$, so that group
is trivial.

If every object has at most two minimal rows, apply
Theorem~\ref{C:categorical-two-row-act-cups}.
The only remaining possibility has four total rows, three total
columns, and row counts $(3,1)$ at objects $(a,b)$.
Indeed five rows would already cost seventeen arrows, and
four columns together with four rows would cost eighteen.
Nonzero second homology forces column counts $(2,1)$ at
$(a,b)$, by the objectwise balanced-tensor presentation.
The band and its two outside identities cost fourteen arrows;
there is at most one extra arrow in $S$.

If every arrow $a\to b$ is minimal, the minimal idempotent
$e_b$ fixes every incoming arrow. There is no external entrance
at $b$. Lemma~\ref{C:one-sided-corner-deletion} deletes $b$
by an actual homotopy equivalence. The resulting ordinary act
cone still has at most twenty arrows and still satisfies all
four conditions, contradicting
Theorem~\ref{C:finite-target-small-act-twenty}.

Otherwise the sole extra internal arrow goes from $a$ to $b$.
The corner at $a$ is exactly the identity-adjoined $3\times2$
band. In the projective augmentation sequence
$0\to V\to P\to\mathbf1\to0$ over $\mathbb QS$, the row
augmentation module has dimension two at $a$ and zero at $b$.
Every nonidentity corner element acts by zero on it, and every
cross map is zero. Hence $V=L\oplus L$ for a one-dimensional
module $L$. Put $L_F=\ker(\mathbb QF\to\mathbf1_{\rm right})$.
The natural categorical act-cone Tor formula and dimension
shifting give
\[
 H_4(B\mathcal C(S,F);\mathbb Q)
 \cong\operatorname{Tor}^{\mathbb QS}_3(L_F,\mathbf1)
 \cong\operatorname{Tor}^{\mathbb QS}_2(L_F,L)^{\oplus2}.
\]
Its dimension is even, contradicting the fourth Betti number
one. This exhausts the possibilities. Every use of the ordinary
act theorem followed an actual homotopy equivalence; no finite
CW type was assumed for an intermediate space.
\end{proof}

\begin{corollary}\label{C:strong-pair-projective-endpoint21}
Let $C$ have three objects, with a strongly connected two-object
source component $S$ and a remaining sink corner $N$. Suppose
the trivial right $\mathbb QN$-module is projective.
If $|C|\leq21$, then $BC\not\simeq\mathbb CP^2$.
Dually, a projective trivial left module at a source corner
gives the same conclusion when the strong pair is the sink.
\end{corollary}
\begin{proof}
If $|N|=1$, apply
Theorem~\ref{C:strong-pair-trivial-endpoint21}.
Otherwise suppose $BC\simeq\mathbb CP^2$ and use the actual
extreme-corner quotient of
Proposition~\ref{C:finite-target-extreme-corner}.
It replaces $N$ by a point and each incoming left $N$-set by
its orbit set, producing a right $S$-act cone $Y$ with
\[
 |\bar C|\leq |C|-|N|+1\leq20.
\]
The quotient is onto on fundamental groups and is a rational
cohomology-ring equivalence. It has a homotopy left inverse
$h:Y\to BC$. If $u$ generates $H^2(BC;\mathbb Z)$, put
$v=h^*u$. The integral argument in
Corollary~\ref{C:finite-target-corner-CP2} applies unchanged:
$Y$ is simply connected, $H^2(Y;\mathbb Z)\cong\mathbb Z$,
$v$ is its generator, and $v^2$ is primitive modulo torsion.
The identities $Q^*v=u$ and $Q^*(v^2)=u^2$ prove the last
two assertions by the split surjections in degrees two and four.
For every finite group $H$ and map $z:BH\to Y$, Miller's
theorem makes $hz$ null-homotopic, so $z^*v=0$ integrally.
Finally the minimal local group of $N$ has order at most
$|N|\leq21$, so $65521$ is a good prime for the same quotient.
Thus $Y$ satisfies all four conditions of the preceding
proposition, a contradiction. Taking opposites proves the dual.
\end{proof}
